\chapter{The digits of the probable prime}\label{app:prime}
\addednote{\revC}

These are the \ThDigits{} decimal digits of $q_{25332}$, the denominator of the $25{,}332$nd convergent of $\pi$ (Section~\ref{sec:th1}), in rows of 80 digits and groups of 10. The SHA-256 hash of the digits written without spaces or line breaks is
\begin{center}\texttt{\ThSha}.\end{center}
\noindent It is a probable prime: it passed the strong Baillie--PSW test and 25 Miller--Rabin rounds, and a proof of its primality is in progress. The same digits are on the companion site as a plain text file, one line with no spaces \citep{RhythmSite}. Anyone can rebuild the number from $\pi$'s continued fraction with the script \texttt{c12\_third\_pi.py}.

% generated by compute/c12_third_pi.py -- do not edit
\begin{flushleft}\ttfamily\scriptsize
2759256854 5368344821 6470675444 9619863621 3374710001 5410953258 8150875771 4922843075\\
6468607451 4718994153 1458542475 4389868197 8938385516 3290758043 9277925476 6223489742\\
9373949113 8449506040 6579614121 9961078695 0041028984 4076278917 7427057150 0386126115\\
0217895325 0996733181 7240139065 9037150094 7041729980 9163350607 8869395877 0625000590\\
1708799103 6112571362 9040756532 1309385501 2452818308 4079421178 3114350677 4002310401\\
9206798935 7145440912 1328637722 1954188647 8199619166 0808031767 3381687809 9797374850\\
6442838213 6498005428 6431487795 1852843322 0125178020 5594132541 6029788872 3678811596\\
6501444739 2386707147 1513935090 1791988211 5034184103 3037057438 0494932063 9534298373\\
8905476908 6363357487 1004018149 4700171961 7005612928 2542341286 1931132337 7330847163\\
8762470684 1471236347 8766722915 6863492516 2118468604 6179626560 9286014046 2382389278\\
9604008836 7736960797 5496445899 6330879438 5042250815 3892106036 0421769039 4850309337\\
5520690352 7433771935 1690479524 5934433689 3698991169 9064494886 0329436509 8237107348\\
2173844891 5183469838 7634874291 3999144222 0714657869 9877545633 9740121403 1804694362\\
1706609658 6720186459 1949834510 8637333357 5175532585 6607747266 9510761471 9756528566\\
1004749603 9322064679 5160044600 6245525768 2211990616 9047124687 8881922505 5440485936\\
7944852066 3492352683 4932655969 0560843324 6743702622 9929736120 0896493351 7804042931\\
3870088938 3445263177 1189105765 2676079868 4428227801 5341081118 7572995746 8710933669\\
5190428570 4169704817 7128095999 2760603088 9337040918 6934436055 4792403153 9783774746\\
4959154063 9597032884 2934891583 1731970676 7034148709 4010339499 3974793466 3554558191\\
1648239820 8579156302 0833443365 8817634601 0796457793 9315257340 7657162007 4801951981\\
8262780979 4499240873 1406424812 1260901490 2794588741 7918035623 8227568087 7431518294\\
6255215563 9646983996 6602317316 2165153080 6211092417 9627740721 2193742151 9794003769\\
7123657851 3400379150 2885978740 7006571330 7281085592 8801282121 3617731982 4446205603\\
9245276314 5808355460 1529541235 4625695638 8439927457 7290673931 1568232783 6194545295\\
3191733899 8387670459 9249187700 7106503392 5160499136 4953137413 2241429021 5298216967\\
8656390829 5095915909 2320016761 2974207673 1740435043 4145514244 5673938790 3195678819\\
7243840343 8084519933 7049616407 0637539295 1860163769 7503193132 7983271430 7922347622\\
3642963110 1359811050 3625526097 6005879133 6763136401 7977731962 9984302354 9933853297\\
3287728670 1386181825 0662701691 9355934052 1355461423 6663258076 3878051781 5492739386\\
4224857519 5540786410 0550404581 2914420021 7446603282 6291761424 7472817074 7143675355\\
8467165482 9946724079 5690689328 8139522350 8432809518 1982402575 4935413378 5519985385\\
3142271394 8590866152 4580477227 4883261577 2636359994 0699508870 7172521637 4343090138\\
8699966999 5059707477 3125366642 2751515104 7762278444 5336462551 0225013819 4900025467\\
3947968471 6332390718 8870568505 7046186208 8539155360 9157322364 0555584546 3940113974\\
8731326299 3145087957 1005243455 1420491768 3192028605 5627183409 0595072059 0689881130\\
1374575055 6474314558 1188825817 1303304661 0862439589 3819546380 4938920324 8948705371\\
4392954707 4212621629 3187192572 9094769282 7442065837 0754824535 0448740552 2267070063\\
2013090532 6435816322 4130731643 9568390971 5887023050 0156635116 6805243327 5612179124\\
2846777181 3222820626 5149005965 6514800888 3882884824 7152352031 5106522935 8530070105\\
9741923582 4088381163 5609904532 3829185930 9340772908 1536928095 4427872192 8372096137\\
7035735320 2567429627 2189813877 6722716485 5387672930 3043783259 8495211326 8450841509\\
5550865807 7259886451 9045768601 2823107086 2521716611 0323880199 8988228012 1962588367\\
9770981384 0283474621 7315445766 0712467737 9497986579 7855225855 2894108845 6298865167\\
4519316064 7644768770 0660352377 1564197531 1925337605 6961433967 6128541824 4757266537\\
6396709730 0555950302 0581893251 9943546589 9014451888 5429359490 3896060984 8612281708\\
0237560353 8760433606 5717494084 9900737195 5168538556 5853454476 1008988299 6002770681\\
2195881453 2467605676 1415108131 5778388118 5471134602 9612095608 5678794343 4721285632\\
2828762502 6727500616 3169579964 7858345733 8595454434 5395658327 6000343615 3466486246\\
8930135084 9617346543 7677009404 9278797404 6947573638 4140226043 8733994785 7259904442\\
1705736172 0947948057 2205322278 9712576750 4574253717 0530156302 3689287686 2317272618\\
4814827192 1430158026 8104718405 9490417080 5713590033 4334288385 1594035967 6812590424\\
8050778612 9975673520 5439050281 1540351234 6724599391 8577227518 3435944753 2485859265\\
8517198920 7093703538 3687528222 9123991624 2534475829 2334008067 1353620180 8511281271\\
4053914835 1437053110 4823532743 5962044276 9035028065 1463700646 0300070581 3352723340\\
2397527373 0265795357 0150872772 0660141820 1861523319 9219024090 6162303006 0872834193\\
9672344262 9658774454 6695267985 4645779183 4025160872 8703581074 2337964682 4575200311\\
1697281237 1758228422 9495811052 4315019346 0748528401 5263648019 7025720628 8632114446\\
1612351594 6389062742 4080615225 2815481140 3220313518 1054449653 0550036723 2774897837\\
4663700840 9013009214 7385220371 7982705413 3246792015 6372985203 9780433496 6470066447\\
3969172847 8415353546 0224987549 8306164384 0833420410 1743176326 0497522435 6393853565\\
5605480636 4140794750 1407731382 3256572399 0997981282 0359414645 4632760431 9575841852\\
3710784311 4610255195 4731660382 6439401765 3710027609 7879771504 5699472582 9557789559\\
4603464890 1962601167 6711369863 8098527680 8539862319 2983362662 9831493025 7349389044\\
9456246549 4326477815 6769828574 8641426246 7135223050 3390574579 5338768615 4034527517\\
4231370355 2106110982 2065876556 1233007683 6842062320 9934182891 7018501019 0160135096\\
3288676016 9703807092 1193880183 2925793671 0671324916 9825700678 1345772892 2763006533\\
3330859967 4412827400 5237555321 7841983829 6271480286 2932391905 2516808091 6026858698\\
9801808505 1014026492 7145206850 1625940566 0110523313 9868386526 9386542339 3478991187\\
5256274905 0971797934 4301086952 5202451033 1269155542 9154016183 4691438440 8059787033\\
8532148986 1661647605 7745111516 4341121362 7632832936 9271313374 5776466592 7739489487\\
6919859144 9375893706 3996298644 8405772914 2432945595 4108537916 2374781638 2713622064\\
0235555145 2839443120 7428885842 9808401353 9277634613 7659750090 3338378792 5995895228\\
3146493722 4596928221 3741562954 3347043083 7273575310 0878687643 7153122576 0927013674\\
9738253284 2377096132 6984031460 5016873261 9485018425 5205156813 4426337115 6249956214\\
1478478473 2190503105 6917450136 7119451723 5848911647 1311281300 5551740838 0008611382\\
1005417140 3628364822 9442902463 4353510628 2738739032 3888387628 9877955866 9532327253\\
1276191667 5585977770 0693026992 5997091098 7954822399 9608328685 5012177302 6561616956\\
7121279022 9557039989 3903292861 4630147373 5146314195 0198279125 0287594348 0123151283\\
7527176943 5044629534 7670343968 6337246485 8547460041 9841343503 6801483157 7130528285\\
9743807672 5428699318 2944753632 9459333365 0283949505 6807096523 1848224665 2000816996\\
3601568841 8722874426 5652354956 5846313123 0720936705 3864981019 4287311028 3481444416\\
7830337423 7732699571 2032557946 9035282654 5671103232 6328052861 9481993205 8882053181\\
7540809994 0763797991 0873418888 8925393802 3116891706 8481886462 1809910942 7471521199\\
2615366034 6326437168 8967125089 0916253398 4513864850 4881543418 4082824655 9892860310\\
6329946132 3585601616 2627499709 4228934906 8596117156 7889277889 8908995512 0664895680\\
2513807008 2200398164 8693884104 2922230080 6436663831 0270343116 3933374221 2430457675\\
6317494995 4441723036 4356170628 4798658861 8904021761 5249454138 1763791457 2420461342\\
3424256442 6962877053 8218425262 9874527744 8136011163 1784916337 0341510242 5812475682\\
5376415636 4785686875 0656248771 9752002820 6023095388 0933234502 6720858946 3872179507\\
1595553565 8594870420 0220632856 4246254540 3281764208 7913944597 6932155154 7465794608\\
9418216431 1258609970 6343371298 5043926796 5778815127 9417014092 4798097072 4505701391\\
1812595347 2170035270 0677945911 9780303861 1798188621 7236141849 6225370349 0992488860\\
5645030619 9720829104 8296733680 5491109370 6967259936 1852037425 0044683471 2197875816\\
7095231102 9855000954 0176547144 5249173971 3126954790 8287233815 5266495150 7533746229\\
4358491312 7282432886 4797935465 6433062890 8344113319 1454582236 1876918754 1977025755\\
9380769199 5447351393 1319095634 9525201933 2996005572 6666987625 6928490798 4356616629\\
5456513092 3584798816 9398968267 7573635086 3842261127 2052020349 5431306366 3510772647\\
5516152985 3598564546 3260718753 2849328144 1766925722 4099888419 2549549021 2441910109\\
5965454768 2420207143 1241583636 7214546205 6634653555 0539891213 1369693173 1671052412\\
9602691502 0924361603 2526745749 9920411840 8071251466 4774740579 6198925583 4402609083\\
5397859504 5777353090 5361821231 1569671875 4319007575 0925677306 2090001696 9736481660\\
5068551571 8735737176 6422442187 1274020073 9265305832 5312821770 9236178019 7305607910\\
7975033141 2863768179 9111837661 2536337351 1804996519 1291406655 6178519840 4436494810\\
3593519946 1287439395 6755881043 6921143449 0724933880 2636746641 3295532287 6533875789\\
7427865021 5171150578 1754873760 8497036684 2471026368 7082417993 3335158102 2979299095\\
4350510042 9820595916 3412503564 6042541637 4814121234 7374027769 2682432611 5527140703\\
3042957088 7126290896 2117782788 7267271104 2565497494 9360687836 2795392942 8018141467\\
8591963515 6616178719 5150839731 0838635542 2755386381 5268411827 1682698535 9649980653\\
7643306587 4286530265 1810257372 9219741371 2796811664 6593335806 0821883683 5826905358\\
8115610749 8586475725 9160443204 1387204386 9320620453 3102718524 4767782652 1000020351\\
9995137616 8427955161 6207012065 4099442721 2425887017 6296318274 0028141046 4653560759\\
4503027850 4362399285 4106603880 2755300491 1680981149 9342950822 6595736593 6816557978\\
9330306917 0176795418 8765065341 2610250226 2898193749 1788188037 4038198345 4721594013\\
3434347974 1267590445 3164517326 2414126261 5634017438 1719236343 3931223822 1799950627\\
1141817201 3387192843 1232282354 6957720171 3204860663 8052181200 6124530371 9336242457\\
3100484292 9407844526 6451096958 2125026382 5218093936 3877482174 4599791928 2669110087\\
4251739320 6531045580 1172573578 2024420985 7977101487 7943275489 4953141688 4703579443\\
9467624043 4229961565 4321993690 3893422792 0592737408 9271222991 4912603108 0384866082\\
7422901756 4326833420 0210079863 3083489035 6010773498 1701268557 4714183938 2779686593\\
5177287732 6590334801 3444313459 8353570705 5807650455 6713621961 3909836113 9178024884\\
2141785544 1815455555 7882396744 8132026732 1841999314 8339499170 3732314061 2204591097\\
8969585294 6806755577 2799652548 1506613175 5187238722 4912464266 0584588980 1145980648\\
1174216950 8090655061 5724661901 9450012043 1701605794 3756483060 0066497894 8650843474\\
9831172120 4472895875 4376958282 1435872420 4452435522 4847430874 6429782298 3840427426\\
1465015039 8124707401 6739548786 3645611208 6836667853 2401182145 6049397985 2899520904\\
3659460131 8039567242 9683741834 6023003246 9693243911 6639207827 5061142360 3668685234\\
6857186667 3158266104 4653284017 3160340739 8556943376 6400783561 3757055608 1888961042\\
2071993105 5877164006 0628295906 3474283106 6992486829 5058682286 9346502047 6892748443\\
1279053929 8375728318 3910265121 7625011475 3423618825 8134110476 8796593362 6033245005\\
7332434477 1503107565 3708371101 7925139043 4413689002 7786742961 6314309176 4994370566\\
1914898865 0715953878 8552980223 8795226063 3553140240 9738503052 0423910851 9478251234\\
8579469687 0652010295 2028586193 5824694907 9190642653 6799658436 8773030116 3843392584\\
8500291370 5520024743 6296015860 2189941969 9299420387 0783192862 4502766505 3907809291\\
2656691237 1669449469 6522675868 5791167250 8661882249 1522934980 9599909629 9395696564\\
9979958608 7718326350 1644152503 1464699164 9637534862 2236556470 3481473692 6434557390\\
3699363818 9202933253 0249894691 9741987555 3221763661 6958869367 7115935187 2777120206\\
6277169857 6645530039 2438034985 8479115125 2806713797 4846527205 2974322155 1024151848\\
4624218883 1218478566 3775189123 4664906617 6122960409 7999351156 7830494514 1448439901\\
1784046296 6579713921 3426729413 2548624132 7513849872 9966072873 3054904265 8931715540\\
5842926737 3378549790 4109413954 1606660476 2632649687 3595494280 6773157460 9142314602\\
0548669001 6986214762 2044446513 9214753383 8694288507 3722090907 3332076045 7349674568\\
0576119095 8328418231 9520628350 2308184050 4555209566 8402752217 8987854008 6333002457\\
1685843888 7610218186 0128624819 1231865757 3661427192 8680664912 9817121580 8008835056\\
7768733271 6785043102 1000639630 9066895833 2929674195 1787660314 1437595533 0783532300\\
5164899954 4920308802 4086580558 7673265943 0519843337 1821514232 8638234355 5460403705\\
6101407880 3846289156 7529062339 7188962085 5258400609 9834916276 0336310387 9620658027\\
2625478249 2651850009 2079811443 6002739273 8314222866 2538550999 6040453109 1764033760\\
8161009355 2419042870 0902524510 4822656876 6236487740 0638027184 0363614023 0795462156\\
7359159623 4470000618 2540192737 3978022184 2237824485 2484639084 4134831838 9196545075\\
8995941426 3354143263 9574440757 0238564911 8185426476 9298350213 4842110250 6015730706\\
6556510072 9259091059 6579888565 3375855937 9116964079 3033372769 6532808262 5289416221\\
7892897730 4049374932 6184428912 0836456585 3282968347 4900330186 3768107727 6220077423\\
7090139627 3990047681 8729704428 3794223217 5032747461 1305963152 2780382774 8500509514\\
0043577400 0510200310 3759518026 6327587714 1702845634 7893295385 4501234909 8034659316\\
9686437819 9422633309 8760288512 4312568098 1747076951 1931576906 0654519846 9794500610\\
6496465494 0120833330 5364394624 8276179666 3531881950 0941154864 0960944432 4708177578\\
3769442474 7675381661 6135957198 6036657989 8157603467 2923778807 2364101624 5863044260\\
9089360278 1110515593 7408778924 7667661220 4650801641 7135769939 2179273840 5873878348\\
8077599135 0356659519 8056720202 4731640765 3122720327 5733258275 3236429505 3181525988\\
3436893485 2630871468 7187936915 2519715432 6589550878 6089250046 4149697596 1123563181\\
9272256275 1490612434 0935119856 8476247826 9581711328 1126368438 0728047500 4720007284\\
1297231039 1537741874 1602586728 2739625340 1134127311 8877287050 9515876812 1883823239\\
6790868102 2927614806 8835468606 9110504103 6014942809 0949703
\end{flushleft}

